\section{Coupled Authorization Failures}
\label{sec:covariance}

Zero-floor enforcement experiments do not reveal how controls fail. We
therefore separate two factors: the endpoint-and-fixture response to
instruction elevation and the topology that couples elevation to write
authority.

\subsection{Risk identity}

Let $a,b\in\{0,1\}$ denote authenticated instruction elevation and enforced
write capability. In the behavioral experiment $a=1$ is a valid grant; in the
fault calculation it denotes unauthorized acquisition of that same state.
Let $Y$ be an exact protected-write request generated while the tool is
visible, with $\rzero=\Pr(Y=1\mid a=0)$ and
$\rone=\Pr(Y=1\mid a=1)$. The gate here is a capability check, so the monitor
mutates the protected resource exactly when $Y=b=1$; because the receiver is
never entitled to it, that mutation is unauthorized (Table~\ref{tab:contract}).
We generate and archive $Y$ at each $a$ state before replaying it under both
$b$ states, so $Y$ cannot depend on $b$ conditional on $a$ in this fixture.
For matched marginals $\E[a]=\E[b]=q$,
\begin{equation}
R=q[(1-q)\rzero+q\rone]+\Cov(a,b)(\rone-\rzero),
\label{eq:identity}
\end{equation}
so relative to independent lifecycles,
\begin{equation}
R_{\mathrm{top}}-R_{\mathrm{ind}}=\Cov(a,b)(\rone-\rzero).
\label{eq:excess}
\end{equation}
This follows by expanding the two $b=1$ cells. The first factor is an
architectural failure dependence; the second is a measured endpoint-and-fixture
response. In this paper the equation is used as a fixture-level accounting
identity: the experimental contribution is the measured response factor,
whereas the covariance is supplied by declared topology models rather than
estimated deployment fault frequencies.
When $\rone\approx\rzero$, channel typing does not change behavior and
capability enforcement remains necessary. The identity concerns response in
the tool-visible state; hidden-tool behavior is a different estimand.

The response is not a property of the endpoint alone. Measured on different
tool-visible fixtures, it was .875 for DeepSeek semantic placement, .396 for
Qwen on the same task (supplement, App.~D), .313 for Qwen credential gating
(\S\ref{sec:credential-gated}), .167 for Qwen exact requests at $a=0$ versus
$a=1$ on 48 fixed relay outputs (\S\ref{sec:e4-gateway}), and .063 in the tool-visible arm of the
capability-hidden fixture (42/48 to 45/48, \S\ref{sec:defense}). The spread
within Qwen is 6.3-fold (.396/.063), whereas the cross-endpoint contrast on a
shared fixture is 2.2-fold (.875/.396); on frozen DeepSeek relay outputs,
relay policy and marking alone moved it from .125 to .771
(Table~\ref{tab:e4-gateway}). Fixture variation therefore exceeds the endpoint
variation we observe, and the design cannot attribute $\rone-\rzero$ to the
model. Figure~\ref{fig:covariance-risk} shows the consequence: under one
declared topology, the computed excess scales with whichever fixture supplied
the response. Tool hiding changes the available action schema, so the hidden-tool
route cannot be inserted into Equation~\eqref{eq:identity}. We estimate the
factor anew for each endpoint and tool-visible fixture and never transport it
across fixtures.

\subsection{Credential-gated measurement}
\label{sec:credential-gated}

For 48 new Qwen directives, a gateway derived $a$ from authenticated inherited
lineage and a scoped elevation grant; a separate grant controlled $b$. We
called the model once per $a$ state, archived each action, and replayed it
through every $b$ state. Exact protected requests increased from 33/48 to
48/48, so $\rone-\rzero=.3125$ (paired 95\% CI [.1866,.4625]). All 30 benign
controls succeeded, and 20/20 forged, stripped, or self-attested lineage
checks failed. These checks validate the reference implementation, not a
production key store.

\begin{table}[t]
  \centering
  \footnotesize
  \caption{Qwen credential-gate risk computed at declared $q=.05$. Only
  $\rzero=33/48$ and $\rone=48/48$ are model measurements; covariance and
  risk are computed from the declared fault topology.}
  \label{tab:measured}
  \setlength{\tabcolsep}{3pt}
  \begin{tabular}{lrrr}
    \toprule
    Topology & $\E[a]{=}\E[b]$ & $\Cov(a,b)$ & Risk (\%) \\
    \midrule
    Independent lifecycles & .05 & 0 & 3.52 \\
    Nominal split, common root & .05 & .0231 & 4.24 \\
    Shared persistent state & .05 & .0475 & 5.00 \\
    \bottomrule
  \end{tabular}
\end{table}

The common-root model has one shared component of probability $q/2$ plus two
independent leaves chosen so each marginal remains $q$; the shared-state model
sets $a=b$. Their computed excesses over independence are 0.72 and 1.48
percentage points, which equal $\Cov(a,b)\times.3125$ to six decimals
(.023141 and .047500 before rounding). Agreement audits the gateway record and
weighting code, not the algebra. The value $q=.05$ is a
sensitivity parameter, not an estimated compromise frequency.

\begin{figure}[t]
  \centering
  % Figure 4: excess expected mutation risk, Cov(a,b) * (r1 - r0), for every
% tool-visible endpoint-and-fixture response measured in the paper. Only the
% slopes are measurements; the covariances are declared topologies at q=.05.
\definecolor{fig4qwen}{RGB}{68,114,196}
\definecolor{fig4deepseek}{RGB}{204,102,0}
\begin{tikzpicture}[font=\footnotesize]
  \begin{axis}[
    scale only axis, width=46mm, height=36mm,
    xmin=0, xmax=0.0475, ymin=0, ymax=4.4,
    axis lines*=left, tick align=outside, clip=false,
    tick label style={font=\scriptsize}, label style={font=\scriptsize},
    xtick={0,0.023141,0.0475},
    xticklabels={0, .0231, .0475},
    scaled x ticks=false,
    ytick={0,1,2,3,4},
    xlabel={$\Cov(a,b)$ (declared, $q=.05$)},
    ylabel={Excess risk (pp)}, ylabel style={yshift=-3pt},
    ymajorgrids, grid style={black!12},
  ]
    % Declared topologies.
    \draw[densely dashed, black!35] (axis cs:0.023141,0) --
      (axis cs:0.023141,4.4);
    \node[font=\tiny, black!60, anchor=south, align=center]
      at (axis cs:0.023141,4.4) {common\\[-2pt]root};
    \node[font=\tiny, black!60, anchor=south, align=center]
      at (axis cs:0.0475,4.4) {shared\\[-2pt]state};
    \node[font=\tiny, black!60, anchor=south, align=center]
      at (axis cs:0,4.4) {indep.};
    % Excess = Cov * (r1 - r0); endpoints at Cov = .0475.
    \addplot[thick, fig4deepseek, mark=*, mark size=1.1pt]
      coordinates {(0,0) (0.023141,2.0248) (0.0475,4.1563)};
    \addplot[thick, fig4qwen, mark=*, mark size=1.1pt]
      coordinates {(0,0) (0.023141,0.9160) (0.0475,1.8802)};
    \addplot[very thick, black, mark=square*, mark size=1.3pt]
      coordinates {(0,0) (0.023141,0.7232) (0.0475,1.4844)};
    \addplot[thick, fig4qwen, densely dashed, mark=*, mark size=1.1pt,
      mark options={solid}]
      coordinates {(0,0) (0.023141,0.3857) (0.0475,0.7917)};
    \addplot[thick, fig4qwen, densely dotted, mark=*, mark size=1.1pt,
      mark options={solid}]
      coordinates {(0,0) (0.023141,0.1446) (0.0475,0.2969)};
    \node[font=\tiny, anchor=west, align=left] at (axis cs:0.0475,4.1563)
      {.875 DeepSeek, semantic};
    \node[font=\tiny, anchor=west, align=left] at (axis cs:0.0475,1.98)
      {.396 Qwen, semantic};
    \node[font=\tiny, anchor=west, align=left] at (axis cs:0.0475,1.40)
      {\textbf{.313 Qwen, credential}};
    \node[font=\tiny, anchor=west, align=left] at (axis cs:0.0475,0.7917)
      {.167 Qwen, frozen relay};
    \node[font=\tiny, anchor=west, align=left] at (axis cs:0.0475,0.2969)
      {.063 Qwen, tool-visible E12};
  \end{axis}
\end{tikzpicture}

  \caption{Declared-fault sensitivity. Points are measured responses
  $\rone-\rzero$ with the write tool visible; bars are 95\% intervals where
  the run report gives one for the difference (E3, \S\ref{sec:credential-gated};
  DeepSeek, where $\rzero=0$). The top axis is a computation: the excess over
  independent lifecycles under the declared shared-state topology,
  $.0475(\rone-\rzero)$ at $q=.05$, not an observed failure rate.}
  \Description{A dot plot of five measured elevation responses: DeepSeek
  semantic placement .875 with interval .748 to .953; Qwen semantic
  placement .396; Qwen credential gate .313 with interval .187 to .463; Qwen
  relay outputs .167; Qwen tool-visible arm of the capability-hidden
  fixture .063. A top axis converts each response to the declared
  shared-state excess, from 0 to 4.75 percentage points.}
  \label{fig:covariance-risk}
\end{figure}

\subsection{Relay policy changes the response factor}
\label{sec:e4-gateway}

\begin{table}[t]
  \centering
  \footnotesize
  \caption{Exact protected requests over frozen DeepSeek relay outputs
  (n=48 per arm). The provider is evaluated at both elevation states.}
  \label{tab:e4-gateway}
  \setlength{\tabcolsep}{3pt}
  \begin{tabular}{lrr}
    \toprule
    Relay policy and source & $a=0$ & $a=1$ \\
    \midrule
    Benign, marked & 0/48 & 6/48 \\
    Adaptive, marked & 1/48 & 32/48 \\
    Benign, unmarked & 0/48 & 37/48 \\
    \bottomrule
  \end{tabular}
\end{table}

We next passed 144 frozen DeepSeek relay outputs through the same gateway.
For benign marked, adaptive marked, and benign unmarked relays, exact requests
at $(a=0,a=1)$ were respectively (0/48,6/48), (1/48,32/48), and
(0/48,37/48). The prespecified adaptive-minus-benign interaction was .5208
($p=4.17\times10^{-7}$). Adaptive and unmarked exact endpoints were not
distinguished ($p=.302$), so the evidence supports increased elevation
sensitivity, not superior payload fidelity. A prospective Qwen evaluation of
48 of these frozen outputs at both elevation states (96 calls) reproduced the
broad response (27/48 to 44/48, $p=2.21\times10^{-4}$) and a
smaller, marginal exact response (22/48 to 30/48, $p=.0386$); a same-bank
repeat gave 25/48 to 43/48 and 21/48 to 30/48. These runs bound one temporal
repeat, not provider-wide variance.

The result is a measured decomposition, not three executed provider arms:
agent calls estimate $\rzero,\rone$; the declared topology weights reconstruct
risk. Separate token names do not guarantee $\Cov(a,b)=0$ when credentials
share an issuer, storage root, principal, or recovery path.
